\section{Point-in-Polygon Testing}

Given a polygon $\mathcal{P}$ and query point $q$, we want to determine
whether $q$ lies inside the polygon. The implementation in this project uses
the winding-number algorithm.

\subsection{Winding Number}

Imagine tracing the polygon boundary while observing how many times the
polygon winds around the query point. For a simple polygon,
\[
WN(q)=
\begin{cases}
0, & \text{if } q \text{ is outside},\\
+1 \text{ or } -1, & \text{if } q \text{ is inside}.
\end{cases}
\]

The sign depends on whether the polygon vertices are ordered
counterclockwise or clockwise. The implementation therefore uses
\[
WN(q)\neq 0
\]
as the inside condition.

\subsection{Upward Crossings}

For an edge from $A$ to $B$, an upward crossing occurs when
\[
y_A\leq y_q \qquad \text{and} \qquad y_B>y_q.
\]
If additionally
\[
\operatorname{orientation}(A,B,q)>0,
\]
the winding number is incremented:
\[
WN\leftarrow WN+1.
\]

\subsection{Downward Crossings}

A downward crossing occurs when
\[
y_B\leq y_q \qquad \text{and} \qquad y_A>y_q.
\]
If
\[
\operatorname{orientation}(A,B,q)<0,
\]
then
\[
WN\leftarrow WN-1.
\]

The asymmetric inequalities prevent polygon vertices lying exactly on the
horizontal query line from being counted twice.

\subsection{Boundary Convention}

For collision detection, this project treats points lying on the polygon
boundary as inside the obstacle. Therefore, before updating the winding
number, each edge is checked using the point-on-segment predicate.

\subsection{Complexity}

Each polygon edge is examined exactly once. For a polygon containing $n$
vertices,
\[
T(n)=O(n).
\]
